% R5: introduction and dataset (Sanika). Builds on the proposal's problem statement.

\section{Introduction and Dataset}\label{sec:intro}

Electric motor systems used 53\% of the world's electricity in 2023, 72\% in industry
\cite{iea2025}. Vibration monitoring catches bearing and rotor faults before they cause a
failure, but the standard diagnostics demodulate bearing-defect frequencies from accelerometers
sampled at kilohertz rates \cite{randall2011}. The common benchmark, the Case Western Reserve
bearing data, is sampled at 12 and 48\,kHz, and standard envelope analysis already diagnoses
many of its faults \cite{smith2015}. Small workshops and farms skip that equipment and run motors to
failure. A smartphone is already at every machine, but its inertial sensor logs at about 100\,Hz.

At 100\,Hz nothing above the 50\,Hz Nyquist frequency survives \cite{shannon1949}: it is removed
or folded back into the band. For the motor studied here (shaft speed $f_1 \approx 15$, 20 and 25\,Hz), the
outer-race defect frequency of a cracked bearing (44--86\,Hz) and the twice-line-frequency
signature of voltage imbalance (60--100\,Hz) lie above that limit, while the rotor-bar sidebands
at $f_1 \pm 2sf_\mathrm{supply}$ \cite{nandi2005} and the impulsiveness caused by poor lubrication
lie below it.
The usual question, which classifier is best, therefore turns into a different one (Track 2:
dimensionality reduction against noise and out-of-distribution data): \emph{which class
structure survives the truncation, which projection preserves it, and how does it hold up when
the sensor degrades or the operating condition changes?} Answering it needs an honest test.
Overlapping windows cut from one recording are near-copies, and random splits over windows
overstate vibration-diagnosis accuracy \cite{wheat2024}, a case of the leakage that inflates
machine-learning results across fields \cite{kapoor2023}, so every result here holds out whole
recordings and is repeated over five seeds. We contribute (i) a separability boundary: which
fault information a 100\,Hz sensor keeps, which it hides, and why; (ii) the size of the leakage
gap, from 1.00 macro-F1 on random splits to 0.47 on unseen recordings; and (iii) how the
projections hold up under sensor corruption, an unseen speed or load, and domain adaptation.
Sections~\ref{sec:methods}--\ref{sec:conclusion} give the methods, results and conclusions.

\textbf{Related work.} Classical pipelines pair hand-crafted time- and frequency-domain features
with PCA or LDA and a shallow classifier \cite{duda2001}; Kernel PCA \cite{scholkopf1998} is the
non-linear variant, Subspace Alignment \cite{fernando2013} adapts such a subspace to a new
operating condition, and DTW \cite{sakoe1978} compares waveforms without features. We test each
at 100\,Hz on held-out recordings.

\textbf{Dataset.} We use the public (CC BY 4.0) dataset of Ertar{\u{g}}{\i}n et al.~\cite{ertargin2026}. An iPhone 15 Pro Max lies on
a squirrel-cage induction motor and logs six inertial channels at 100\,Hz: the raw accelerations
$g_X, g_Y, g_Z$ and the gravity-compensated user accelerations $g^u_X, g^u_Y, g^u_Z$. There are
six health states: healthy (H), insufficient lubrication (B1), severely insufficient lubrication
(B2), a cracked outer bearing ring (B3), supply-voltage imbalance (V) and a broken rotor bar (R).
Each class is recorded once at each of six operating conditions (supply frequency 30, 40 or
50\,Hz, unloaded or loaded), giving 36 balanced recordings of about 15 minutes (90{,}000 samples
per channel).

\textbf{Pre-processing.} Each recording is cut into windows of 1{,}024 samples (10.24\,s, a
0.098\,Hz frequency resolution) with a 512-sample step: 6{,}264 labelled windows, 174 per
recording. The mean of each window is removed per channel, which drops gravity and sensor
offset. Amplitudes are not normalised, because absolute vibration energy is discriminative;
every scaling step is fitted inside the cross-validation folds.
